\documentclass[border=10pt]{standalone}
\usepackage{tikz,ctex}
\usetikzlibrary{positioning}

\begin{document}

\begin{tikzpicture}[
    % 全局设置
    node distance = 1.6cm and 1.4cm,
    every node/.style = {
        draw,
        rounded corners,
        thick,
        align = center,
        font = \sffamily\small,
        minimum width = 2.0cm,
        minimum height = 0.9cm
    },
    % 样式定义
    centerStyle/.style = {
        fill = blue!45,
        text = white,
        font = \sffamily\bfseries\large,
        line width = 2pt,
        minimum size = 2.5cm,
        circle
    },
    branchTitle/.style = {
        font = \sffamily\bfseries,
        line width = 1.5pt
    },
    subNode/.style = {
        fill = white,
        font = \sffamily\small,
        minimum width = 2.0cm,
        minimum height = 0.8cm
    },
    leafNode/.style = {
        fill = white,
        font = \sffamily\footnotesize,
        minimum width = 1.8cm,
        minimum height = 0.7cm
    },
    arrowStyle/.style = {
        ->,
        thick,
        >=stealth,
        shorten >=2pt
    }
]

    % ================= 1. 中心节点 =================
    \node[centerStyle] (center) {梯度计算 \\ Evaluation of Gradients};

    % ================= 2. 上方分支：反向传播基础 (红) =================
    \node[branchTitle, fill=red!15, above=of center] (bp) {反向传播基础 \\ Backprop Basics};

    \node[subNode, above left=of bp] (single) {单层网络 \\ Single-Layer};
    \node[subNode, above=of bp] (general) {一般前馈网络 \\ General Feed-Forward};
    \node[subNode, above right=of bp] (delta) {误差项 $\delta_j$ \\ Error Terms};
    \node[subNode, right=of bp, yshift=0cm] (autodiff) {自动微分 \\ Automatic Differentiation};

    % ================= 3. 右侧分支：雅可比矩阵 (蓝) =================
    \node[branchTitle, fill=orange!15, right=of center] (jacobian) {雅可比矩阵 \\ Jacobian Matrix};

    \node[subNode, above=of jacobian] (sens) {局部敏感度 \\ Local Sensitivity};
    \node[subNode, above right=of jacobian] (perturb) {输入扰动传播 \\ Input Perturbation};
    \node[subNode, right=of jacobian] (recursive) {递归反向传播 \\ Recursive Backprop};
    \node[subNode, below right=of jacobian, yshift=0cm] (forward) {前向传播法 \\ Forward Propagation};
    \node[subNode, below=of jacobian] (numdiff) {数值微分检验 \\ Numerical Check};

    % ================= 4. 下方分支：黑塞矩阵 (绿) =================
    \node[branchTitle, fill=green!15, below=of center] (hessian) {黑塞矩阵 \\ Hessian Matrix};

    \node[subNode, below left=of hessian] (second) {二阶导数 \\ Second Derivatives};
    \node[subNode, below=of hessian] (complexity) {计算复杂度 \\ Complexity};
    \node[subNode, below right=of hessian] (diag) {对角近似 \\ Diagonal Approximation};
    \node[subNode, left=of hessian, yshift=0cm] (outer) {外积近似 \\ Outer Product};
    \node[subNode, right=of hessian] (lm) {Levenberg--Marquardt \\ 近似};

    % ================= 5. 左侧分支：数值微分 (紫) =================
    \node[branchTitle, fill=purple!15, left=of center] (numeric) {数值微分 \\ Numerical Differentiation};

    \node[subNode, above left=of numeric] (finite) {有限差分 \\ Finite Differences};
    \node[subNode, left=of numeric] (central) {中心差分 \\ Central Differences};
    \node[subNode, below left=of numeric] (cost) {计算成本 $\mathcal{O}(W^2)$ \\ Computational Cost};
    \node[subNode, above=of numeric, yshift=-0.3cm] (verify) {正确性检验 \\ Correctness Check};

    % ================= 连线逻辑 =================
    % 中心 -> 四大分支标题
    \draw[arrowStyle, red] (center) -- (bp);
    \draw[arrowStyle, blue] (center) -- (jacobian);
    \draw[arrowStyle, green!60!black] (center) -- (hessian);
    \draw[arrowStyle, purple] (center) -- (numeric);

    % 分支标题 -> 子节点 (上方)
    \draw[arrowStyle, red!60] (bp) -- (single);
    \draw[arrowStyle, red!60] (bp) -- (general);
    \draw[arrowStyle, red!60] (bp) -- (delta);
    \draw[arrowStyle, red!60] (bp) -- (autodiff);

    % 分支标题 -> 子节点 (右侧)
    \draw[arrowStyle, blue!60] (jacobian) -- (sens);
    \draw[arrowStyle, blue!60] (jacobian) -- (perturb);
    \draw[arrowStyle, blue!60] (jacobian) -- (recursive);
    \draw[arrowStyle, blue!60] (jacobian) -- (forward);
    \draw[arrowStyle, blue!60] (jacobian) -- (numdiff);

    % 分支标题 -> 子节点 (下方)
    \draw[arrowStyle, green!60!black] (hessian) -- (second);
    \draw[arrowStyle, green!60!black] (hessian) -- (complexity);
    \draw[arrowStyle, green!60!black] (hessian) -- (diag);
    \draw[arrowStyle, green!60!black] (hessian) -- (outer);
    \draw[arrowStyle, green!60!black] (hessian) -- (lm);

    % 分支标题 -> 子节点 (左侧)
    \draw[arrowStyle, purple!60] (numeric) -- (finite);
    \draw[arrowStyle, purple!60] (numeric) -- (central);
    \draw[arrowStyle, purple!60] (numeric) -- (cost);
    \draw[arrowStyle, purple!60] (numeric) -- (verify);

\end{tikzpicture}

\end{document}